% Source: Corsin Nick/Class Notes/Week 10.pdf, pp. 3--12
\chapter{Week 10 --- Geodesics, Flux \& the Divergence Theorem}
\label{ch:week10}

% Transition: Claude Sonnet 5
This chapter opens with a short, informal derivation of what \qt{shortest path} means on a
submanifold, using the length functional from Week 9 as an optimization target. It then turns
to \newterm{flux}: the operational fix for a lecture-notes definition of a bounded $C^1$
domain that Corsin finds unusable directly, via interior/exterior normals and the Gauss map,
and Gauss' divergence theorem, which converts a flux integral over $\partial\Omega$ into a
volume integral over $\Omega$. The exercise sheet is at the end of the chapter, in
\Cref{sec:week10_solutions}.

\begin{ainote}
Corsin opens this week's notes by restating, verbatim, the definitions of $L(\gamma)$ and
$L(c)$ from \cref{def:length_c1_curve,def:length_c0_curve} in Week 9 --- the same page,
reproduced as a recap before the geodesics discussion below, which uses the $C^1$ length
functional as its starting point. Not repeated here; see Week 9.
\end{ainote}

\section{A short introduction to geodesics}
\label{sec:geodesics}

\begin{ainote}
Corsin marks this section \qt{not exam relevant}. It is included anyway, since the calculation
is short and the underlying idea --- treat length as a functional and apply optimization to it
--- is the same one used later in the course for genuinely examinable variational arguments.
\end{ainote}

We will try to find the shortest path from $0 \in \mathbb{R}^n$ to $x \in \mathbb{R}^n$. This
is, of course, the straight line, but the essentially same calculation can be used to find the
shortest path between points on a sphere or any other submanifold.

\textbf{Idea.} Apply \qt{optimization} to the map
\[L : \mathcal A \to [0,\infty), \qquad c \mapsto L(c) = \int_0^1 \lvert c'(t)\rvert\,dt,
  \qquad \mathcal A := \{\gamma \in C^1([0,1],\mathbb{R}^n) : \gamma(0)=0,\ \gamma(1)=x\}.\]

\begin{lemma}[A few differentiation identities]
\label{lem:inner_product_derivative_identities}
Let $\alpha,\beta : (0,1) \to \mathbb{R}^n$ be $C^1$. Then
\[\frac{d}{dt}\langle\alpha(t),\beta(t)\rangle = \langle\alpha'(t),\beta(t)\rangle
  + \langle\alpha(t),\beta'(t)\rangle, \qquad
  \frac{d}{dt}\lvert\alpha(t)\rvert = \frac{1}{\lvert\alpha(t)\rvert}\langle\alpha(t),\alpha'(t)\rangle.\]
\end{lemma}

% Source: Corsin Nick/Class Notes/Week 10.pdf, pp. 4--5
\begin{proposition}[The shortest path is a straight line]
\label{prop:geodesic_straight_line}
Assume that $c \in C^\infty([0,1],\mathbb{R}^n)$ with $c(0)=0$, $c(1)=x$ and
$\lvert c'(t)\rvert \equiv \lambda$ is the optimal path we are looking for. Then $c(t) = xt$
and $\lambda = \lvert x\rvert$.
\end{proposition}

\begin{proof}
Let $g \in C^\infty([0,1],\mathbb{R}^n)$ with $g(0)=g(1)=0$, and define, for $\varepsilon \in
\mathbb{R}$, the family of paths $c_\varepsilon(t) := c(t) + \varepsilon g(t)$, so that
$c_0 = c$.

% Sonnet 5 (Medium) --- FIG-W10-02
\begin{center}
\begin{tikzpicture}[>=Latex, scale=1.2]
  \draw[red, thick] (0,0) -- (3.6,1.1) node[midway, below, font=\scriptsize, sloped] {$c$};
  \foreach \k in {1,2,3} {
    \draw[green!45!black, thick, opacity={0.5+0.15*\k}]
      plot[smooth] coordinates {(0,0) (1.1,{0.55+0.13*\k}) (2.3,{0.75+0.10*\k}) (3.6,1.1)};
  }
  \node[font=\scriptsize] at (2.0,1.55) {$c_\varepsilon$};
  \fill (0,0) circle (1.3pt) node[left, font=\scriptsize] {$0$};
  \fill (3.6,1.1) circle (1.3pt) node[right, font=\scriptsize] {$x$};
\end{tikzpicture}
\end{center}

By assumption, $c$ is a minimizer of $L$, so it is a critical point of the family, i.e.
\[\left.\frac{d}{d\varepsilon}\right|_{\varepsilon=0} L(c_\varepsilon) = 0.\]
(This equality itself requires a proof, but the idea is exactly the same as that of
optimization in $\mathbb{R}^n$.) First,
\[\left.\frac{\partial}{\partial\varepsilon}\right|_{\varepsilon=0} c_\varepsilon(t)
  = \left.\frac{\partial}{\partial\varepsilon}\right|_{\varepsilon=0}
    \big(c(t)+\varepsilon g(t)\big) = g(t).\]
Recalling $\lvert c'\rvert \equiv \lambda$ and using
\Cref{lem:inner_product_derivative_identities},
\[
\begin{aligned}
\left.\frac{d}{d\varepsilon}\right|_{\varepsilon=0} L(c_\varepsilon)
  &= \left.\frac{d}{d\varepsilon}\int_0^1 \lvert c_\varepsilon'(t)\rvert\,dt\right|_{\varepsilon=0}
   = \int_0^1 \left.\frac{\partial}{\partial\varepsilon}\lvert c_\varepsilon'(t)\rvert\,dt
     \right|_{\varepsilon=0} \\
  &= \int_0^1 \frac{\big\langle c_\varepsilon'(t),\ \partial_\varepsilon c_\varepsilon'(t)
       \big\rangle}{\lvert c_\varepsilon'(t)\rvert}\,dt\bigg|_{\varepsilon=0}
   = \frac1\lambda \int_0^1 \Big\langle c'(t),\ \partial_\varepsilon\big|_{\varepsilon=0}
       c_\varepsilon'(t)\Big\rangle\,dt \\
  &= \frac1\lambda \int_0^1 \frac{d}{dt}\Big\langle c'(t), g(t)\Big\rangle
       - \big\langle c''(t), g(t)\big\rangle\,dt
   = \frac1\lambda\left[\big\langle c'(t),g(t)\big\rangle\Big|_0^1
       - \int_0^1 \big\langle c''(t),g(t)\big\rangle\,dt\right] \\
  &= -\frac1\lambda \int_0^1 \big\langle c''(t),g(t)\big\rangle\,dt \ \overset{!}{=}\ 0,
\end{aligned}
\]
where the boundary term vanished since $g(0)=g(1)=0$. Since $g$ is arbitrary, it follows that
$c''(t) \equiv 0$, and together with $c(0)=0$, $c(1)=x$ this gives $c(t)=xt$ and $\lambda =
\lvert x\rvert$.
\end{proof}

\begin{ainote}
Corsin flags a subtlety in his own proof with a warning triangle at the step
$\left.\frac{d}{d\varepsilon}\right|_{\varepsilon=0}L(c_\varepsilon) = 0$ and again at the final
$= 0$: both are exactly the two places where genuine work is being hidden. The first
is the claim that a minimizer of $L$ over the infinite-dimensional space $\mathcal A$ is a
critical point of every one-parameter family through it --- the same fact used for
$\mathbb{R}^n$-optimization in Week 5, now needing $\mathcal A$ to be a vector space (it is
not: $\mathcal A$ is an affine subspace, but $c_\varepsilon - c \in \{g : g(0)=g(1)=0\}$, a genuine
vector space, which is what makes the argument go through). The second is the
\newterm{fundamental lemma of the calculus of variations}: that
$\int_0^1\langle c''(t),g(t)\rangle\,dt = 0$ for \emph{every} admissible $g$ forces
$c''\equiv0$ pointwise, not just on average. Corsin does not name or prove this lemma; it is
the standard fact that if a continuous function integrates to zero against every test function
vanishing at the endpoints, it is identically zero (take $g$ a bump concentrated where $c''$
is supposedly nonzero).
\end{ainote}

\section{Bounded \texorpdfstring{$C^1$}{C1} domains and why the definition in the script is so terrible}
\label{sec:bounded_c1_domains}

% Source: Corsin Nick/Class Notes/Week 10.pdf, p. 6
% Def. 14.3
\begin{definition}[Bounded $C^k$ domain]
\label{def:bounded_ck_domain}
Given $k \geq 1$, a bounded open set $\Omega \subset \mathbb{R}^n$ is called a
\newterm{bounded $C^k$ domain} if $M := \partial\Omega$ is an $(n-1)$-dimensional submanifold
of class $C^k$.
\end{definition}

Citing the 2024 lecture notes, Corsin builds up the machinery needed to actually \emph{use}
this definition --- to define \newterm{flux} we need the concept of interior and exterior
normal vectors. For $\Omega$ as above, a normal vector $v \in N_p(\partial\Omega)$ is
\newterm{interior} if
\[p + hv \in \Omega \quad \text{for all } h>0 \text{ small},\]
and \newterm{exterior} if
\[p + hv \in \mathbb{R}^n\setminus\Omega \quad \text{for all } h>0 \text{ small}.\]

A \newterm{Gauss map} is a continuous map of unit normal vectors on $\partial\Omega$. There
exists a unique interior and exterior Gauss map; write
\[\nu \in C^0(\partial\Omega, \mathbb{S}^{n-1}) \text{ such that } \nu(p) \text{ is exterior
  normal}\]
for the exterior unit normal map.

\begin{theorem}[Jordan--Brouwer separation theorem]
\label{thm:jordan_brouwer}
If $M \subseteq \mathbb{R}^{n+1}$ is an $n$-dimensional compact and connected submanifold,
then $\mathbb{R}^{n+1}\setminus M$ has exactly two connected components, an
\newterm{inside} and an \newterm{outside}.
\end{theorem}

\section{Flux and the divergence theorem}
\label{sec:flux_divergence_theorem}

% Source: Corsin Nick/Class Notes/Week 10.pdf, p. 7
\begin{definition}[Flux]
\label{def:flux}
Suppose $F : U \to \mathbb{R}^n$ is a $C^1$ vector field, $U \subseteq \mathbb{R}^n$ open, and
$\Omega \subseteq U$ is a $C^1$ bounded domain. The \newterm{flux} (\germanterm{Fluss}) of $F$
through $\partial\Omega$ is defined as
\[\int_{\partial\Omega} \langle F,\nu\rangle\,d\vol_{n-1} \qquad \text{(analysis notation)}
  \quad = \quad \int_{\partial\Omega} \vec F(x)\cdot d\vec A \qquad \text{(physics notation)}\]
\[= \int_B \big\langle F\circ\phi(x),\ \nu\circ\phi(x)\big\rangle
   \sqrt{\det D\phi(x)\transp D\phi(x)}\,dx \qquad \text{(explicit)},\]
where $\nu : \partial\Omega \to \mathbb{S}^{n-1}$ is the exterior unit normal, i.e.\ $\nu(y)$
is the outward-facing unit normal of $\partial\Omega$ at $y \in \partial\Omega$, and
$\phi : B \to \partial\Omega$ is a (local) parametrization of $\partial\Omega$. In general,
multiple integrals in the explicit formulation may be necessary.
\end{definition}

% Source: Corsin Nick/Class Notes/Week 10.pdf, p. 8
\begin{proposition}[Flux in $\mathbb{R}^3$]
\label{prop:flux_r3}
Let $\Omega \subseteq \mathbb{R}^3$ be a bounded $C^1$ domain and $F \in
C^1(\overline\Omega,\mathbb{R}^3)$ a vector field, and let $\phi : B \to \partial\Omega
\setminus N$ be a parametrization of $\partial\Omega\setminus N$, where $N \subseteq
\partial\Omega$ is a null set, $B \subseteq \mathbb{R}^2$. Then:
\begin{enumerate}[label=\textbf{\arabic*.}]
  \item $\displaystyle\nu(x) := \frac{\partial_1\phi(x)\times\partial_2\phi(x)}
    {\lvert\partial_1\phi(x)\times\partial_2\phi(x)\rvert}$ is an (interior or exterior) Gauss
    map (more precisely, $\nu\circ\phi^{-1}$ is a Gauss map, i.e.\
    $\nu(x) \in N_{\phi(x)}(\partial\Omega\setminus N)$).
  \item The Gram determinant is given by
    $\sqrt{\det D\phi(x)\transp D\phi(x)} = \lvert\partial_1\phi(x)\times\partial_2\phi(x)
    \rvert$.
  \item Therefore,
    \[\int_{\partial\Omega} \vec F\cdot d\vec A = \pm\int_B \big\langle F\circ\phi(x),\
      \partial_1\phi(x)\times\partial_2\phi(x)\big\rangle\,dx,\]
    where the sign depends on whether $\nu(x)$ is interior ($-$) or exterior ($+$).
\end{enumerate}
\end{proposition}

% Source: Corsin Nick/Class Notes/Week 10.pdf, pp. 9--10
\begin{example}[Flux through the unit sphere]
\label{ex:flux_unit_sphere}
Compute the flux of $F(x,y,z) := (xz,\,z,\,y)$ through the surface of $\mathbb{S}^2$ with
interior $B_1(0)$.

We use spherical coordinates with $r\equiv1$:
\[\phi : (0,2\pi)\times(0,\pi) \to \mathbb{R}^3, \qquad (\varphi,\theta) \mapsto
  \begin{pmatrix}\sin\theta\cos\varphi \\ \sin\theta\sin\varphi \\ \cos\theta\end{pmatrix}.\]

\textbf{Normal vector.}
\[\partial_\varphi\phi = \begin{pmatrix}-\sin\theta\sin\varphi\\ \sin\theta\cos\varphi\\
  0\end{pmatrix}, \qquad \partial_\theta\phi = \begin{pmatrix}\cos\theta\cos\varphi\\
  \cos\theta\sin\varphi\\ -\sin\theta\end{pmatrix}
  \quad \Longrightarrow \quad \partial_\varphi\phi\times\partial_\theta\phi
  = \begin{pmatrix}-\sin^2\theta\cos\varphi\\ -\sin^2\theta\sin\varphi\\
  -\sin\theta\cos\theta\end{pmatrix}.\]

\textbf{Sanity check: does this make sense, and is it exterior or interior?} We have
\[\partial_\varphi\phi\times\partial_\theta\phi = -\sin\theta\ \phi(\varphi,\theta)
  = -\sin\theta\,(x,y,z),\]
and the exterior normal of $\mathbb{S}^2$ is $+(x,y,z)$, so $\partial_\varphi\phi\times
\partial_\theta\phi$ is an \emph{interior} normal (for $\theta\in(0,\pi)$, $\sin\theta>0$).
So we multiply by $-1$ to switch the order in the cross product, matching the exterior sign
in \cref{prop:flux_r3}.

\textbf{The integral.}
\[
\begin{aligned}
\int_{\mathbb{S}^2} \vec F\cdot d\vec A
  &= \int_0^{2\pi}\!\!\int_0^\pi \begin{pmatrix}xz\\ z\\ y\end{pmatrix}\cdot
     \begin{pmatrix}\sin\theta\cos\varphi\\ \sin\theta\sin\varphi\\ \cos\theta\end{pmatrix}
     \sin\theta\,d\theta\,d\varphi \\
  &= \int_0^{2\pi}\!\!\int_0^\pi \Big[\sin^3\theta\cos\theta\cos^2\varphi
     + \sin^2\theta\cos\theta\sin\varphi + \sin^2\theta\cos\theta\sin\varphi\Big]\,d\theta\,
     d\varphi \\
  &= \pi\int_0^\pi \sin^3\theta\cos\theta\,d\theta
   = \pi\left[\frac{\sin^4\theta}{4}\right]_0^\pi = 0,
\end{aligned}
\]
using $\int_0^{2\pi}\cos^2\varphi\,d\varphi = \pi$ and $\int_0^{2\pi}\sin\varphi\,d\varphi=0$
for the middle two terms.
\end{example}

% Source: Corsin Nick/Class Notes/Week 10.pdf, p. 11
\begin{theorem}[Gauss]
\label{thm:gauss_divergence}
Let $\Omega \subseteq \mathbb{R}^n$ be a bounded $C^1$ domain and $F : \overline\Omega \to
\mathbb{R}^n$ a $C^1$ vector field. Then
\[\int_\Omega \divg F\,d\vol_n = \int_{\partial\Omega} \langle F,\nu\rangle\,d\vol_{n-1},
  \qquad \text{or, in physics notation,} \qquad
  \int_\Omega \vec\nabla\cdot\vec F\,dx = \int_{\partial\Omega} \vec F\cdot d\vec A.\]
\end{theorem}

\begin{ainote}
Corsin dates the theorem to Lagrange (discovered 1762) and Ostrogradsky (proven 1828), and
adds, next to the physics form, an unusually enthusiastic \qt{how cool is this!} --- a volume
integral of a derivative reduced to a boundary integral of the original field, the same shape
as the fundamental theorem of calculus.
\end{ainote}

\begin{example}[Recomputing \cref{ex:flux_unit_sphere} via the divergence theorem]
\label{ex:flux_unit_sphere_via_divergence}
With $\Omega = B_1(0)$ and $F(x,y,z) = (xz,z,y)$ as before, $\divg F = \partial_x(xz) +
\partial_y(z) + \partial_z(y) = z + 0 + 0 = z$. So, in spherical coordinates ($z =
r\cos\theta$),
\[
\int_{\partial\Omega} \vec F\cdot d\vec A = \int_\Omega \divg F\,dx = \int_\Omega z\,dx
  = \int_0^1\!\!\int_0^{2\pi}\!\!\int_0^\pi \underbrace{\cos\theta\sin\theta}_{=\,\frac12\sin(2\theta),\
    \pi\text{-periodic}} r^3\,d\theta\,d\varphi\,dr = 0,
\]
since $\int_0^\pi \sin(2\theta)\,d\theta = 0$ --- confirming \cref{ex:flux_unit_sphere} with
none of its cross-product bookkeeping.
\end{example}

\begin{ainote}
Corsin writes the divergence as $\divg F = z$ directly without the intermediate
$\partial_x(xz)+\partial_z(z)+\partial_y(y)$ step; spelled out above for clarity. Worth
noting how much the two computations of the same flux differ in character: the direct
definition needs an explicit Gauss map, an orientation sanity check, and three cross terms
in the integrand, while the divergence theorem needs only $\divg F$ and a coordinate change
already used in \cref{ex:paraboloid_patch_area} of Week 9. This is the theorem's main
practical use --- not conceptual elegance, but that $\divg F$ is often far simpler than $F$
restricted to a curved boundary.
\end{ainote}

% Source: Corsin Nick/Class Notes/Week 10.pdf, p. 12
Gauss' theorem is also a powerful tool in physics. For example, the volume $V$ of water in a
moving fluid with velocity $\vec v$ passing through a surface element $d\vec A$ in time
$\Delta t$ is
\[V = \Delta t\,\vec v\cdot d\vec A\cdot\varrho \qquad (\varrho \text{ the density}),\]
since a fluid column of length $L = \Delta t\,\lvert\vec v\rvert$ and cross-section $d\vec A$
tilted at angle $\theta$ to $\vec v$ sweeps a volume $L\lvert d\vec A\rvert\cos\theta =
\Delta t\,\vec v\cdot d\vec A$ through the surface element.

% Sonnet 5 (Medium) --- FIG-W10-03
\begin{center}
\begin{tikzpicture}[>=Latex, scale=1.1]
  \draw[purple, thick] (0,0) -- (2.4,0) -- (3.1,1.3) -- (0.7,1.3) -- cycle;
  \draw[red, thick] (2.4,0) -- (2.4,1.5);
  \foreach \y in {0.15,0.45,0.75,1.05} {
    \draw[orange, thick, ->] (-1.1,\y) -- (0.9+0.29*\y,\y);
  }
  \draw[orange, thick, ->] (1.55,0.65) -- (2.85,1.15) node[right, font=\scriptsize] {$\vec v$};
  \draw[green!50!black] (2.2,0.55) arc (160:70:0.5);
  \node[green!50!black, font=\scriptsize] at (2.15,0.85) {$\theta$};
  \node[font=\scriptsize] at (-1.3,0.6) {$\vec v$};
  \node[purple, font=\scriptsize] at (1.55,-0.25) {$L=\Delta t\lvert\vec v\rvert$};
  \node[red, font=\scriptsize, right] at (2.55,0.75) {$h=\lvert d\vec A\rvert\cos\theta$};
\end{tikzpicture}
\end{center}

Therefore, the fluid flowing out of any bounded $C^1$ domain $\Omega$ per unit time is
\[\int_{\partial\Omega} \varrho\vec v\cdot d\vec A = \int_\Omega \vec\nabla\cdot(\varrho
  \vec v)\,dx.\]
Of course, since the total amount of fluid is conserved, this is also the change in the
amount of fluid in $\Omega$ over time:
\[\int_\Omega \vec\nabla\cdot(\varrho\vec v)\,dx = -\frac{d}{dt}\int_\Omega \varrho\,dx
  = -\int_\Omega \frac{\partial\varrho}{\partial t}\,dx.\]
Since this holds for any $\Omega$, the \newterm{continuity equation}
(\germanterm{Kontinuit\"atsgleichung}) follows:
\[\vec\nabla\cdot(\varrho\vec v) + \frac{\partial\varrho}{\partial t} = 0.\]

\exercisesheet{10}

\textit{Statements quoted verbatim from \texttt{exercises/Ex10\_Analysis2\_eng.pdf} (assigned
24 April 2026, due 4 May 2026). Attribution: Prof.\ Joaquim Serra, D-MATH, ETH Z\"urich.}

% Source: Corsin Nick/Class Notes/Week 10.pdf, p. 1
\begin{ainote}[Corsin's recommendations]
Colour key, as given on the page: blue \textbf{= important}, orange \textbf{= semi-important},
red \textbf{= optional}.

\begin{center}
\begin{tabular}{@{}l l p{0.5\linewidth}@{}}
\textbf{Problem} & \textbf{Priority} & \textbf{Corsin's note} \\
\hline
10.1 & important & ``Use spherical coordinates with $r\equiv1$.'' \\[3pt]
10.2 & important & (worked parametrization of the torus --- see below) \\[3pt]
10.3 & important & ``1.)--4.), 5.)'' \\[3pt]
10.4 & semi-important & --- \\[3pt]
10.5 & semi-important & --- \\[3pt]
10.6 & semi-important & --- \\
\end{tabular}
\end{center}
\end{ainote}

\begin{ainote}
Only the exercises tagged important above are reproduced below: 10.1--10.3. 10.4--10.6
(tractrix, logarithmic spiral, moment of inertia of an ellipse) are not reproduced.
\end{ainote}

% Source: Corsin Nick/Class Notes/Week 10.pdf, p. 1
\begin{ainote}
For \textbf{10.3}, Corsin's note \qt{1.)--4.), 5.)} appears to flag that the starred
sub-part \textbf{5.} --- the harder \qt{if and only if} characterization of a bounded $C^1$ domain,
marked (*) on the official sheet --- is worth attempting alongside the more routine parts
1--4, not skipped as the (*) marker might suggest.
\end{ainote}

% Source: exercises/Ex10_Analysis2_eng.pdf, p. 1
\begin{exercise}[10.1 --- Spherical quadrilateral \textnormal{(important)}]
\label{ex:10.1}
Calculate the area of the spherical quadrilateral $S \subset \mathbb{S}^2$ bounded by the
meridians $-\pi < \varphi_1 < \varphi_2 < \pi$ and parallels $-\pi/2 < \theta_1 < \theta_2 <
\pi/2$.
\end{exercise}

% Source: exercises/Ex10_Analysis2_eng.pdf, p. 1
\begin{exercise}[10.2 --- The surface area of a torus \textnormal{(important)}]
\label{ex:10.2}
Let $S := \{(x,y,0)\in\mathbb{R}^3 \mid x^2+y^2=4\}$ and $T := \{x\in\mathbb{R}^3 \mid d_S(x)
\leq 1\}$, where $d_S(x) = \inf\{\lvert x-y\rvert \mid y\in S\}$ denotes the minimal distance
from $S$ to $x\in\mathbb{R}^3$. Parametrize $\partial T$ and then calculate the area.

% Source: Corsin Nick/Class Notes/Week 10.pdf, p. 1
\textit{Corsin's hint, general parametrization of a torus surface:}
\[x = (R+r\cos\phi)\cos\theta, \qquad y = (R+r\cos\phi)\sin\theta, \qquad z = r\sin\phi,\]
for the torus swept by a tube of radius $r$ at distance $R$ from the axis.

% Sonnet 5 (Medium) --- FIG-W10-01
\begin{center}
\begin{tikzpicture}[>=Latex, scale=1.0]
  \begin{scope}
    \draw[->] (0,0) -- (2.2,0) node[right, font=\scriptsize] {$x$};
    \draw[->] (0,0) -- (0,1.8) node[above, font=\scriptsize] {$z$};
    \draw[->] (0,0) -- (-1.1,-0.7) node[left, font=\scriptsize] {$y$};
    \draw[blue!70!black, thick] (0.4,0.55) circle (0.9 and 0.35);
    \draw[blue!70!black, thick] (0.05,0.15) circle (1.5 and 0.6);
    \draw[blue!70!black, dashed] (-1.0,0.0) arc (180:360:1.5 and 0.6);
  \end{scope}
  \begin{scope}[xshift=5.2cm]
    \draw[blue!70!black, thick] (0,0) circle (1.6);
    \draw[blue!70!black, thick] (0,0) circle (0.7);
    \draw[->] (0,0) -- (55:1.15) node[midway, above, font=\scriptsize] {$R$};
    \draw[orange, ->] (0,0) -- (18:0.7) node[midway, below, font=\scriptsize] {};
    \draw[orange, thick, ->] (18:0.7) -- (55:1.6)
      node[midway, above right, font=\scriptsize] {$r$};
    \draw (0.55,0) arc (0:18:0.55);
    \node[font=\scriptsize] at (0.75,0.15) {$\theta$};
    \draw[->] (1.1,0) -- (2.0,0) node[right, font=\scriptsize] {$x$};
    \draw[->] (1.1,0) -- (1.1,0.9) node[above, font=\scriptsize] {$y$};
  \end{scope}
  \begin{scope}[xshift=10.4cm]
    \draw[purple, dotted] (-1.6,0) -- (1.6,0);
    \draw (0,0) circle (0.9);
    \draw[->] (0,0) -- (35:0.9) node[midway, above right, font=\scriptsize] {$r$};
    \draw (0.35,0) arc (0:35:0.35);
    \node[font=\scriptsize] at (0.55,0.15) {$\varphi$};
    \draw[->] (-1.6,-0.3) -- (-1.6,0.9) node[above, font=\scriptsize] {$z$};
    \draw[->] (-1.6,-0.3) -- (-0.6,-0.3) node[right, font=\scriptsize] {$x$};
  \end{scope}
\end{tikzpicture}
\end{center}
\end{exercise}

% Source: exercises/Ex10_Analysis2_eng.pdf, p. 1
\begin{exercise}[10.3 --- Solids of revolution \textnormal{(important)}]
\label{ex:10.3}
Given $f \in C^1((a,b))\cap C([a,b])$ such that $f\geq0$, define the \textnormal{rotational
body around the $x$-axis} as
\[\begin{aligned}
  \Omega &:= \left\{(x,y,z)\in\mathbb{R}^3 \mid a<x<b,\ \sqrt{y^2+z^2}<f(x)\right\}, \\
  \partial_{\mathrm{side}}\Omega &:= \left\{(x,y,z)\in\mathbb{R}^3 \mid a<x<b,\ \sqrt{y^2+z^2}=f(x)\right\}.
\end{aligned}\]
\begin{enumerate}[label=\textbf{\arabic*.}]
  \item Say whether $\Omega$, $\partial_{\mathrm{side}}\Omega$ are open/closed/connected
    subsets of $\mathbb{R}^3$.
  \item Show that $\vol_3(\Omega) = \pi\int_a^b f(x)^2\,dx$.
  \item Show that $\vol_2(\partial_{\mathrm{side}}\Omega) = 2\pi\int_a^b
    \sqrt{1+f'(x)^2}\,f(x)\,dx$. \textit{Hint:} parametrize $(x,\theta)\mapsto (x,f(x)\cos\theta,
    f(x)\sin\theta)$.
  \item Calculate the volume and surface area of the `improper rotational body' that arises
    when $f(x)=1/x$, $a=1$, $b=\infty$.
  \item \textnormal{(*)} Show that $\Omega$ is a bounded $C^1$ domain (as in
    \Cref{def:bounded_ck_domain}) if and only if
    \begin{equation}
    f(x)>0 \text{ in } (a,b), \qquad f'(a^+)=+\infty, \qquad f'(b^-)=-\infty. \tag{$*$}
    \end{equation}
\end{enumerate}
\end{exercise}

\section{Solutions}
\label{sec:week10_solutions}

\begin{ainote}[TODO --- pending the remaining exercise statements]
Corsin's Week 10 notes contain no worked solutions to the official sheet (only the torus hint
reproduced above), so all solutions here are original. Only \textbf{10.1--10.3} are solved,
matching the exercises reproduced above; 10.4--10.6 are left for a later pass.
\end{ainote}

\begin{exercisesolution}[Solution to \cref{ex:10.1}]
% Generator: Claude Sonnet 5 (Medium)
We parametrize $\mathbb{S}^2$ (up to a null set) using latitude/longitude spherical
coordinates
\[\Phi : (-\pi,\pi)\times\left(-\tfrac\pi2,\tfrac\pi2\right) \to \mathbb{S}^2, \qquad
  (\varphi,\theta) \mapsto \begin{pmatrix}\cos\varphi\cos\theta\\ \sin\varphi\cos\theta\\
  \sin\theta\end{pmatrix}\]
--- note this is \emph{not} the $(r\equiv1)$ parametrization used elsewhere in this chapter,
whose polar angle runs over $(0,\pi)$; here $\theta$ is measured from the equator, matching
the exercise's range $\theta\in(-\pi/2,\pi/2)$. Computing
\[\partial_\varphi\Phi\times\partial_\theta\Phi = \begin{pmatrix}-\sin\varphi\cos\theta\\
  \cos\varphi\cos\theta\\ 0\end{pmatrix} \times \begin{pmatrix}-\cos\varphi\sin\theta\\
  -\sin\varphi\sin\theta\\ \cos\theta\end{pmatrix} = \begin{pmatrix}\cos\varphi\cos^2\theta\\
  \sin\varphi\cos^2\theta\\ \sin\theta\cos\theta\end{pmatrix},\]
its length is $\lvert\partial_\varphi\Phi\times\partial_\theta\Phi\rvert = \cos\theta$ (for
$\theta\in(-\pi/2,\pi/2)$, so $\cos\theta>0$). By \cref{prop:flux_r3}(2), this is the Gram
determinant, so $dA = \cos\theta\,d\theta\,d\varphi$, and
\[\int_S dA = \int_{\varphi_1}^{\varphi_2}\!\!\int_{\theta_1}^{\theta_2} \cos\theta\,d\theta\,
  d\varphi = \boxed{\ (\varphi_2-\varphi_1)\big(\sin\theta_2-\sin\theta_1\big)\ }.\]
\end{exercisesolution}

\begin{exercisesolution}[Solution to \cref{ex:10.2}]
% Generator: Claude Sonnet 5 (Medium)
The circle $S$ has radius $2$, and $T$ is the solid torus swept by a tube of radius $1$
around it, so $\partial T$ is Corsin's general torus surface with $R=2$, $r=1$ fixed. Using
his parametrization,
\[\phi(\theta,\varphi) = \begin{pmatrix}(2+\cos\varphi)\cos\theta\\
  (2+\cos\varphi)\sin\theta\\ \sin\varphi\end{pmatrix}, \qquad \theta,\varphi\in(0,2\pi).\]
Then $\partial_\theta\phi = \big(-(2+\cos\varphi)\sin\theta,\ (2+\cos\varphi)\cos\theta,\
0\big)$ has length $2+\cos\varphi$, and $\partial_\varphi\phi = (-\sin\varphi\cos\theta,\
-\sin\varphi\sin\theta,\ \cos\varphi)$ has length $1$; the two are orthogonal (their dot
product is $\sin\varphi\cos\varphi\sin\theta\cos\theta - \sin\varphi\cos\varphi\sin\theta
\cos\theta = 0$), so
\[\sqrt{\det D\phi\transp D\phi} = \lvert\partial_\theta\phi\rvert\,\lvert\partial_\varphi\phi
  \rvert = 2+\cos\varphi.\]
Hence
\[\vol_2(\partial T) = \int_0^{2\pi}\!\!\int_0^{2\pi} (2+\cos\varphi)\,d\theta\,d\varphi
  = 2\pi\left(2\cdot2\pi + \underbrace{\int_0^{2\pi}\cos\varphi\,d\varphi}_{=0}\right)
  = \boxed{\ 8\pi^2\ }.\]
\end{exercisesolution}

\begin{exercisesolution}[Solution to \cref{ex:10.3}]
% Generator: Claude Sonnet 5 (Medium)
\textbf{Part 1.} $\partial_{\mathrm{side}}\Omega$ is not closed (it omits, e.g., the limit
point $(a,f(a),0)$ as $x\to a^+$) and not open (a $2$-dimensional surface has empty interior
in $\mathbb{R}^3$), but it \emph{is} path-connected: from any point, first move along the
graph of $f$ in the $x$-direction, then move around the circle in the $(y,z)$-variables.
$\Omega$ is open (as the strict sublevel set of a continuous function), not closed, and its
connectedness depends on $f$: if $f$ vanishes somewhere in $(a,b)$, the cross-section pinches
to a point there and $\Omega$ is disconnected into the pieces on either side; if $f>0$
throughout $(a,b)$, $\Omega$ is connected.

\textbf{Part 2.} In cylindrical coordinates $(x,\rho,\theta)$ with $y=\rho\cos\theta$,
$z=\rho\sin\theta$, $dx\,dy\,dz = \rho\,d\rho\,d\theta\,dx$, the domain becomes the box-like
region $\{a<x<b,\ 0<\rho<f(x),\ \theta\in[0,2\pi)\}$, so
\[\vol_3(\Omega) = \int_a^b\!\!\int_0^{2\pi}\!\!\int_0^{f(x)} \rho\,d\rho\,d\theta\,dx
  = \int_a^b\!\!\int_0^{2\pi} \frac12 f(x)^2\,d\theta\,dx
  = \boxed{\ \pi\int_a^b f(x)^2\,dx\ }.\]

\textbf{Part 3.} Using the suggested parametrization $\Phi(x,\theta) = (x,f(x)\cos\theta,
f(x)\sin\theta)$, write the orthonormal frame $\xi_1:=(1,0,0)$, $\xi_2:=(0,\cos\theta,
\sin\theta)$, $\xi_3:=(0,-\sin\theta,\cos\theta)$, so that
\[\partial_x\Phi = \xi_1 + f'(x)\xi_2, \qquad \partial_\theta\Phi = f(x)\xi_3.\]
Since $\xi_1,\xi_2,\xi_3$ are orthonormal and right-handed, $\xi_1\times\xi_3=-\xi_2$ and
$\xi_2\times\xi_3=\xi_1$, so
\[\partial_x\Phi\times\partial_\theta\Phi = f(x)\big(\xi_1\times\xi_3
  + f'(x)\,\xi_2\times\xi_3\big) = f(x)\big({-\xi_2}+f'(x)\xi_1\big), \qquad
  \lvert\partial_x\Phi\times\partial_\theta\Phi\rvert = f(x)\sqrt{1+f'(x)^2}.\]
By \cref{prop:flux_r3}(2), this is the Gram determinant, so
\[\vol_2(\partial_{\mathrm{side}}\Omega) = \int_a^b\!\!\int_0^{2\pi} f(x)\sqrt{1+f'(x)^2}\,
  d\theta\,dx = \boxed{\ 2\pi\int_a^b f(x)\sqrt{1+f'(x)^2}\,dx\ }.\]

\textbf{Part 4.} With $f(x)=1/x$ on $[1,\infty)$, both integrands are non-negative, so the
improper integrals are unambiguous (possibly $+\infty$). The volume is finite:
\[\pi\int_1^\infty \frac{1}{x^2}\,dx = \pi\Big[-\frac1x\Big]_1^\infty = \boxed{\ \pi\ }.\]
For the surface area, $f'(x) = -1/x^2$, so the integrand is $\sqrt{1+x^{-4}}/x \geq 1/x$
pointwise, and since $\int_1^\infty \frac1x\,dx$ diverges, so does
\[2\pi\int_1^\infty \frac{\sqrt{1+x^{-4}}}{x}\,dx = \boxed{\ +\infty\ }.\]
This is the classical \newterm{Gabriel's horn} (\germanterm{Torricellis Trompete}): finite
volume, infinite surface area.

\textbf{Part 5.} ($\Leftarrow$) Suppose $(*)$ holds. Since $f'(a^+)=+\infty$, the inverse
function theorem applies to $f|_{(a,a+\varepsilon)}$ for $\varepsilon$ small, giving a $C^1$,
increasing inverse there; extend it by
\[g(y) := \begin{cases} f^{-1}(y) & f(a)<y<f(a)+\varepsilon, \\ a & y\leq f(a). \end{cases}\]
Since $f'(a^+)=+\infty$, $\lim_{y\to f(a)^+}g'(y) = \lim_{x\to a^+} 1/f'(x) = 0 =
\lim_{y\to f(a)^-}g'(y)$, so $g\in C^1$ across $y=f(a)$. Around a boundary point $P :=
(a,f(a)\cos\phi,f(a)\sin\phi)$, the map $(t,\theta)\mapsto\big(g(t),f(g(t))\cos\theta,
f(g(t))\sin\theta\big)$ is then a regular $C^1$ parametrization of $\partial\Omega$ (using
$f(a)>0$ to keep the angular direction non-degenerate). The symmetric argument at $b$, using
$f'(b^-)=-\infty$, handles the other end, so $\partial\Omega$ is a $C^1$ submanifold and
$\Omega$ is a bounded $C^1$ domain.

($\Rightarrow$) Suppose $\Omega$ is a bounded $C^1$ domain. If $f$ vanished at some interior
point, $\Omega$ would be disconnected by Part 1 --- excluded by \cref{def:bounded_ck_domain},
which requires $\Omega$ itself to make sense as a single bounded open set with $(n-1)$-manifold
boundary (implicit in the exercise's setup); so $f>0$ on $(a,b)$. For the endpoint condition
at $a$: the point $P:=(a,f(a),0)\in\partial\Omega$ is a smooth boundary point, so
$T_P(\partial\Omega)$ is $2$-dimensional and contains
\[\frac{d}{dt}\Big|_{t=0^+}\big(a,f(a)-t,0\big) = -e_2, \qquad
  \frac{d}{dt}\Big|_{t=0^-}\big(a,f(a)\cos t,f(a)\sin t\big) = f(a)e_3,\]
so $T_P(\partial\Omega) = \operatorname{span}\{e_2,e_3\}$. But $T_P(\partial\Omega)$ must
also contain the limit of the (rescaled, if necessary) tangent vectors
\[v_\varepsilon := \frac{d}{dt}\Big|_{t=0^+}\big(a+\varepsilon+t,\ f(a+\varepsilon+t),\
  0\big) = \big(1,\ f'(a+\varepsilon),\ 0\big)\]
as $\varepsilon\to0^+$ (these are tangent to $\partial\Omega$ at nearby points converging to
$P$). If $f'(a+\varepsilon)$ stayed bounded along some sequence $\varepsilon_k\to0^+$, the
direction of $v_{\varepsilon_k}$ would converge to some $\alpha e_1+\beta e_2 \notin
\operatorname{span}\{e_2,e_3\}$ (its $e_1$-component is the nonzero constant $1$, up to
normalization), contradicting $T_P(\partial\Omega)=\operatorname{span}\{e_2,e_3\}$. So
$f'(a^+)=+\infty$; the symmetric argument at $b$ gives $f'(b^-)=-\infty$.
\end{exercisesolution}

\begin{ainote}
All parts of \cref{ex:10.1,ex:10.2,ex:10.3} were worked out independently before consulting
\texttt{Sol10\_Analysis2\_eng.pdf}. The three main answers ($(\varphi_2-\varphi_1)(\sin\theta_2
-\sin\theta_1)$, $8\pi^2$, $\pi\int_a^bf(x)^2dx$ / $2\pi\int_a^bf(x)\sqrt{1+f'(x)^2}dx$ /
volume $\pi$, area $+\infty$) match exactly. Part 5 was reconstructed independently along the
same two-directional strategy (build an explicit $C^1$ parametrization near each endpoint
using the inverse function theorem; for the converse, derive a contradiction from the tangent
space being forced to both contain $\operatorname{span}\{e_2,e_3\}$ and the limiting
direction of $v_\varepsilon$) and agrees with the official solution in substance, though the
write-up above is condensed relative to the official one.
\end{ainote}
